% Created 2026-03-13 Fri 14:20
% Intended LaTeX compiler: pdflatex
\documentclass[11pt]{article}
\usepackage[utf8]{inputenc}
\usepackage[T1]{fontenc}
\usepackage{graphicx}
\usepackage{longtable}
\usepackage{wrapfig}
\usepackage{rotating}
\usepackage[normalem]{ulem}
\usepackage{amsmath}
\usepackage{amssymb}
\usepackage{capt-of}
\usepackage{hyperref}
\author{The Prologos Project - From Zachary Larson}
\date{2026-03-13}
\title{The Categorical Structure of Layered Recovery: Stratified Endofunctors on Quantale-Enriched Lattices}
\hypersetup{
 pdfauthor={The Prologos Project - From Zachary Larson},
 pdftitle={The Categorical Structure of Layered Recovery: Stratified Endofunctors on Quantale-Enriched Lattices},
 pdfkeywords={},
 pdfsubject={},
 pdfcreator={Emacs 30.2 (Org mode 9.7.39)}, 
 pdflang={English}}
\begin{document}

\maketitle
\tableofcontents

\section{Abstract}
\label{sec:org8545c93}

We identify the precise categorical structure underlying the Layered Recovery
Principle as implemented in three subsystems of the Prologos language: the logic
engine (stratified negation), the session runtime (effect ordering), and the type
checker (reactive constraint resolution). Each system faces the same fundamental
problem: recovering non-monotone behavior on a monotone propagator substrate.

We show that all three systems instantiate a common categorical pattern: a
\textbf{stratified endofunctor} on a \textbf{quantale-enriched lattice}, where the
endofunctor decomposes as a chain of monotone functors between lattices connected
by Galois connections, with non-monotone operations isolated at precisely one
barrier in the chain. The global fixpoint is the least fixpoint of the
endofunctor, computed by Kleene iteration over the stratified chain.

We prove that this structure is sufficient to recover any non-monotone operation
on a monotone substrate, provided the non-monotone operation satisfies a
\emph{stratifiability condition}: it must be expressible as a barrier between two
monotone phases. We characterize the class of non-monotone operations that
satisfy this condition and identify variants for operations that do not.
\section{1. The Three Systems}
\label{sec:org6d1645b}

Three subsystems of Prologos independently discovered and implemented the same
architectural pattern. We analyze each in turn, extracting the precise algebraic
structure.
\subsection{1.1 System I: The Logic Engine (\texttt{rel})}
\label{sec:org363ad45}

\emph{Non-monotone operation}: Negation-as-failure (NAF). If P is not derivable at
the current stratum, conclude ¬P. This is non-monotone because adding a fact P
to the database can cause ¬P to become invalid.

\emph{Implementation}: Stratified evaluation with three layers.

\begin{verbatim}
Layer 1 (Propagator Network): Unification constraints propagate through cells.
  Values refine monotonically. The network converges to a fixpoint.

Layer 2 (ATMS): Choice points create assumptions. Adding nogoods (inconsistent
  combinations) shrinks the set of valid worldviews — still monotone (the nogood
  set only grows, the valid worldview set only shrinks).

Layer 3 (Stratification Controller): Evaluate ¬P only after stratum containing
  P has reached fixpoint. This is the non-monotone barrier.
\end{verbatim}

\emph{Lattice structure}: Let U be the unification lattice (substitutions ordered by
refinement). Each stratum Sₖ operates on U. The join is unification; the bottom
is the empty substitution.

\emph{The stratified chain}:

\begin{verbatim}
S₀ ──lfp──→ U₀* ──neg──→ S₁ ──lfp──→ U₁* ──neg──→ S₂ ──lfp──→ ...
 ↑                                                              ↓
 └──────────────────── Kleene iteration ────────────────────────┘
\end{verbatim}

Where:
\begin{itemize}
\item lfp denotes the least fixpoint of monotone propagation within a stratum
\item neg denotes the non-monotone negation evaluation at the barrier
\item Uₖ* is the fixpoint of stratum k
\end{itemize}

\emph{Algebraic structure}:
\begin{itemize}
\item The unification lattice U is a quantale: substitutions compose (sequential
application σ ∘ τ) and this composition distributes over joins (most general
unifiers). The ordering is refinement: σ ≤ τ iff τ is a refinement of σ.
\item Negation evaluation neg : U* → Init(U) is not a quantale morphism (it is
anti-monotone on the truth ordering). It maps fixpoint states to initial
conditions for the next stratum.
\item Each stratum's propagation is an endofunctor Fₖ : U → U that is monotone.
The fixpoint Uₖ* = lfp(Fₖ).
\end{itemize}

\emph{Sources}: Logic Engine Design (\texttt{2026-02-24\_LOGIC\_ENGINE\_DESIGN.org} §5),
Effectful Computation on Propagators (\texttt{EFFECTFUL\_COMPUTATION\_ON\_PROPAGATORS.org} §2).
\subsection{1.2 System II: The Session Runtime (\texttt{proc})}
\label{sec:org090afd4}

\emph{Non-monotone operation}: IO effect execution. Writing to a file, sending a
message, opening a connection — these are irreversible observations that cannot
be undone, retried, or merged.

\emph{Implementation}: Five-layer architecture.

\begin{verbatim}
Layer 1 (Session Advancement): Session cells advance monotonically through
  protocol states. Each advancement is a Lamport clock tick.

Layer 2 (Data-Flow Analysis): Analyze variable bindings to derive cross-channel
  ordering edges. Monotone edge accumulation.

Layer 3 (Transitive Closure): Compute the complete partial order over effects
  by propagating ordering edges to closure. Monotone fixpoint.

Layer 4 (ATMS Branching): Maintain per-branch ordering hypotheses. Adding
  nogoods only shrinks valid worldviews — monotone.

Layer 5 (Effect Handler): Execute effects in a valid linearization of the
  resolved partial order. This is the non-monotone barrier.
\end{verbatim}

\emph{Lattice structure}: Two lattices connected by a Galois connection.

Let S be the \emph{session lattice}: session types ordered by advancement (bot →
concrete → advanced states → end). The join is the most-advanced common state.

Let E be the \emph{effect position lattice}: sets of ordering constraints (i < j)
over effect indices, ordered by inclusion. The join is set union. Transitive
closure is a closure operator on E (monotone, extensive, idempotent).

\emph{The Galois connection (α, γ)}:

\begin{verbatim}
α : S → E     (extract causal position from session state)
γ : E → S     (reconstruct remaining protocol from effect position)

α(s) ≤ e  ⟺  s ≤ γ(e)    (adjunction property)
\end{verbatim}

Both α and γ are monotone. The key insight: α preserves sequential composition.

\emph{Algebraic structure}:
\begin{itemize}
\item Both S and E are \textbf{effect quantales} (Katsumata 2014, Gordon 2021):
\begin{itemize}
\item Lattice structure (partial order with joins)
\item Sequential composition (⊕) distributes over joins
\item Ordering respects composition: i ≤ i' ∧ j ≤ j' ⟹ i ⊕ j ≤ i' ⊕ j'
\end{itemize}
\item The Galois connection (α, γ) is a \textbf{quantale morphism}: it preserves both the
lattice ordering and sequential composition.
\item Effect execution (Layer 5) is non-monotone: it selects a linearization of the
partial order and performs irreversible IO.
\end{itemize}

\emph{The stratified chain}:

\begin{verbatim}
S ──α──→ E ──tc──→ E* ──exec──→ World
↑    Galois     monotone     non-monotone
│   connection   closure       barrier
│
(session advancement propagators — monotone)
\end{verbatim}

Where tc denotes transitive closure (a monotone closure operator).

\emph{Sources}: Session Types as Effect Ordering (\texttt{2026-03-06\_SESSION\_TYPES\_AS\_EFFECT\_ORDERING.org} §4-6),
Effectful Computation (\texttt{EFFECTFUL\_COMPUTATION\_ON\_PROPAGATORS.org} §3-4).
\subsection{1.3 System III: The Type Checker (Stratified Quiescence)}
\label{sec:org9f622f1}

\emph{Non-monotone operation}: Trait resolution commitment. Selecting a specific
instance from the set of candidates is a \emph{choice} — once committed, the choice
is irreversible. Instance selection can be anti-monotone: adding a more specific
instance can invalidate a previous selection.

\emph{Implementation}: Stratified quiescence with three strata.

\begin{verbatim}
Stratum 0 (Type Propagation): Cell writes (meta solutions → type cells),
  unification propagators, cross-domain bridges. Monotone: values only refine
  in the type lattice.

Stratum 1 (Readiness Detection): Constraint-readiness propagators scan cell
  states to derive "constraint C is ready to retry" signals. Monotone: once
  a cell becomes non-bot, readiness is permanent.

Stratum 2 (Resolution Commitment): Consume readiness signals. Execute
  resolution (retry constraint, resolve trait, solve meta). Non-monotone:
  commitment selects one instance, excluding alternatives.
\end{verbatim}

\emph{Lattice structure}: Three lattices forming a chain.

Let T be the \emph{type cell lattice}: metavariable values ordered by refinement
(bot → partial types → ground types). The merge is unification.

Let R be the \emph{readiness lattice}: sets of ready-constraint descriptors, ordered
by inclusion. The join is set union. Readiness is monotone in T: if a cell is
non-bot, the constraint is ready; more refined cells are still non-bot.

Let A be the \emph{action lattice}: multisets of resolution actions (trait
resolution, constraint retry, hasmethod resolution), ordered by inclusion. Each
action is a data descriptor (free monad pattern) — not an executed effect.

\emph{The functors between strata}:

\begin{verbatim}
α : T → R     (extract readiness from type cell states)
β : R → A     (map readiness signals to resolution actions)
γ : A → T     (commit resolutions → new type cell writes)
\end{verbatim}

\begin{itemize}
\item α is monotone: more refined type cells produce (weakly) more readiness
signals.
\item β is monotone: more readiness signals produce more action descriptors.
\item γ is \textbf{not monotone} in the general case: committing a resolution can
invalidate other pending resolutions (e.g., solving meta ?X = Nat might
make a trait constraint solvable but also make another constraint
unsatisfiable). However, γ is monotone \emph{per action} — each individual
resolution writes a more refined value to a type cell.
\end{itemize}

\emph{The stratified chain (one iteration)}:

\begin{verbatim}
T ──lfp(S0)──→ T* ──α──→ R ──lfp(S1)──→ R* ──β──→ A ──commit──→ T'
↑                                                              ↓
└──────────────── Kleene iteration (fuel=100) ────────────────┘
\end{verbatim}

\emph{Algebraic structure}:
\begin{itemize}
\item T is a quantale: type substitutions compose and composition distributes over
the unification join.
\item R is a join-semilattice (readiness sets under union). It lacks the sequential
composition that would make it a full quantale.
\item A is a free monoid (action sequences). The ordering is prefix/subset.
\item α : T* → R is a Galois connection's upper adjoint (abstraction). The lower
adjoint γ : R → T does not exist in the traditional sense because commitment
is non-monotone. Instead, β ∘ γ : A → T is a \emph{partial function} — it succeeds
only when the action is still valid (re-check guard).
\end{itemize}

\emph{Confluence conditions} (proved in Track 2 Phase 9):
\begin{itemize}
\item S0: Confluent by construction (lattice merge properties).
\item S1: Confluent by construction (readiness is monotone in cell state).
\item S2: Confluent under the \emph{non-overlapping instance invariant}: each constraint
has at most one valid resolution. Prologos enforces this — monomorphic
instances are looked up by key, parametric instances use most-specific-wins
with same-specificity ties rejected as ambiguous (HKT-7).
\end{itemize}

\emph{Sources}: Track 2 Design (\texttt{2026-03-13\_TRACK2\_REACTIVE\_RESOLUTION\_DESIGN.md} §3.4),
Track 2 Phase 9 (\texttt{test-resolution-confluence-01.rkt}).
\section{2. The Common Categorical Structure}
\label{sec:org847acb7}

\subsection{2.1 Definition: Stratified Recovery System}
\label{sec:orgfd76091}

A \textbf{stratified recovery system} is a tuple \((L, F_m, f_{nm}, n)\) where:

\begin{enumerate}
\item \textbf{L} is a complete lattice (the \emph{substrate}).
\item \textbf{F\textsubscript{m} : L → L} is a monotone endofunctor (the \emph{propagation phase}).
By Knaster-Tarski, F\textsubscript{m} has a least fixpoint lfp(F\textsubscript{m}).
\item \textbf{f\textsubscript{nm} : L → L} is a (possibly non-monotone) function (the \emph{barrier
operation}).
\item \textbf{n ∈ ℕ ∪ \{ω\}} is the \emph{stratification depth} (number of strata).
\end{enumerate}

The \textbf{stratified fixpoint} is computed by:

\begin{verbatim}
x₀ = ⊥
For k = 0, 1, ..., n-1:
  xₖ* = lfp(F_m | initialized at xₖ)
  xₖ₊₁ = f_nm(xₖ*)
Result = x_n*
\end{verbatim}

This is Kleene iteration on the composition f\textsubscript{nm} ∘ lfp(F\textsubscript{m}), starting from ⊥.
\subsection{2.2 The Enrichment: Quantale Structure}
\label{sec:orgeb940ec}

In all three Prologos systems, the substrate lattice L carries additional
algebraic structure: it is a \textbf{quantale}.

A quantale is a complete lattice (L, ≤, ⊔) equipped with an associative binary
operation (⊗ : L × L → L) called \emph{sequential composition} that distributes over
arbitrary joins:

\begin{verbatim}
a ⊗ (⊔ᵢ bᵢ) = ⊔ᵢ (a ⊗ bᵢ)     (left distributivity)
(⊔ᵢ aᵢ) ⊗ b = ⊔ᵢ (aᵢ ⊗ b)     (right distributivity)
\end{verbatim}

The three systems' quantale structures:

\begin{center}
\begin{tabular}{llll}
System & Lattice L & Join (⊔) & Composition (⊗)\\
\hline
Logic & Substitutions & Most general unifier & Substitution application\\
Sessions & Session states & Most-advanced state & Protocol continuation\\
Types & Type cells & Unification & Type substitution composition\\
\end{tabular}
\end{center}
\subsection{2.3 The Decomposition: Inter-Stratum Galois Connections}
\label{sec:org6fbdade}

The stratified fixpoint computation decomposes the endofunctor
(f\textsubscript{nm} ∘ lfp(F\textsubscript{m})) through intermediate lattices connected by Galois connections.

In the general case, a stratified recovery system with k strata has
intermediate lattices L₀, L₁, \ldots{}, Lₖ and Galois connections:

\begin{verbatim}
(αᵢ, γᵢ) : Lᵢ ⇌ Lᵢ₊₁     for i = 0, ..., k-2

where αᵢ is monotone (abstraction)
      γᵢ is monotone (concretization)
      αᵢ(x) ≤ y  ⟺  x ≤ γᵢ(y)
\end{verbatim}

The non-monotone operation f\textsubscript{nm} factors through the final lattice Lₖ₋₁:

\begin{verbatim}
f_nm = γ₀ ∘ ... ∘ γₖ₋₂ ∘ barrier ∘ αₖ₋₂ ∘ ... ∘ α₀
\end{verbatim}

where \emph{barrier} : Lₖ₋₁ → Lₖ₋₁ is the non-monotone operation at the top
stratum. All other operations in the chain are monotone.

The three systems instantiate this pattern:

\begin{center}
\begin{tabular}{llll}
System & Lattice chain & Galois connections & Barrier\\
\hline
Logic & U → U & identity (single lattice, iterated) & neg : U* → Init(U)\\
Sessions & S → E → E* & (α,γ) : S ⇌ E (quantale morphism) & exec : E* → World\\
Types & T → R → A & α : T* → R (abstraction) & commit : A → T\\
\end{tabular}
\end{center}
\subsection{2.4 Theorem: Stratification Preserves Fixpoint Equivalence}
\label{sec:org61786d8}

\textbf{Theorem} (Fixpoint Equivalence). Let \((L, F_m, f_{nm}, n)\) be a stratified
recovery system. If f\textsubscript{nm} is \emph{stratifiable} — i.e., the composition
f\textsubscript{nm} ∘ lfp(F\textsubscript{m}) is monotone — then the stratified fixpoint equals the direct
fixpoint:

\begin{verbatim}
lfp(f_nm ∘ lfp(F_m)) = stratified_fixpoint(L, F_m, f_nm, n)
\end{verbatim}

\emph{Proof sketch}: Define G = f\textsubscript{nm} ∘ lfp(F\textsubscript{m}) : L → L. Since lfp(F\textsubscript{m}) is
monotone (as a function of the initialization point, by the parametric fixpoint
theorem — Cousot \& Cousot 1979) and f\textsubscript{nm} is stratifiable (monotone when composed
with lfp), G is monotone. By Knaster-Tarski, G has a least fixpoint lfp(G).
The stratified computation is exactly Kleene iteration of G starting from ⊥:

\begin{verbatim}
x₀ = ⊥,  xₖ₊₁ = G(xₖ) = f_nm(lfp(F_m | xₖ))
\end{verbatim}

For monotone G on a complete lattice, Kleene iteration converges to lfp(G)
(possibly at a transfinite ordinal for general complete lattices; at ω for
algebraic/continuous lattices, which all three systems are). □

\textbf{Corollary} (Confluence). If the stratified fixpoint equals the direct fixpoint
and S0 is confluent (lattice merge properties), then the overall computation is
confluent — the result is independent of the order in which propagators fire
within each stratum.

\emph{Note}: The stratifiability condition is not always satisfied. System III (type
checker) satisfies it only under the non-overlapping instance invariant. If
overlapping instances existed, f\textsubscript{nm} ∘ lfp(F\textsubscript{m}) would not be monotone (adding
type information could change which instance is selected, leading to a different
fixpoint). This is why Prologos rejects overlapping instances — it is a
\emph{semantic requirement for confluence}, not merely a design choice.
\subsection{2.5 The Common Structure: Diagram}
\label{sec:org75e2a55}

\begin{verbatim}
                ┌─────────────────────────┐
                │   STRATIFIED RECOVERY    │
                │       SYSTEM             │
                └───────────┬─────────────┘
                            │
          ┌─────────────────┼─────────────────┐
          │                 │                 │
┌─────────▼──────┐ ┌───────▼────────┐ ┌──────▼──────────┐
│ System I: rel  │ │ System II: proc│ │ System III: tc  │
│ (Logic Engine) │ │ (Sessions)     │ │ (Type Checker)  │
└─────────┬──────┘ └───────┬────────┘ └──────┬──────────┘
          │                │                 │
          ▼                ▼                 ▼
┌──────────────────────────────────────────────────────┐
│         QUANTALE-ENRICHED LATTICE  (L, ≤, ⊗)        │
│                                                      │
│  Monotone endofunctor Fm : L → L                     │
│  Fixpoint computation: lfp(Fm) via propagator network│
└──────────────────────┬───────────────────────────────┘
                       │
          ┌────────────▼──────────────┐
          │  GALOIS CONNECTION CHAIN   │
          │                           │
          │  L₀ ⇌ L₁ ⇌ ... ⇌ Lₖ₋₁   │
          │  (αᵢ, γᵢ) monotone pairs  │
          └────────────┬──────────────┘
                       │
          ┌────────────▼──────────────┐
          │  NON-MONOTONE BARRIER     │
          │                           │
          │  barrier : Lₖ₋₁ → L₀     │
          │  (negation / execution /  │
          │   commitment)             │
          └────────────┬──────────────┘
                       │
          ┌────────────▼──────────────┐
          │  KLEENE ITERATION         │
          │                           │
          │  Iterate (barrier ∘ lfp)  │
          │  until global fixpoint    │
          └───────────────────────────┘
\end{verbatim}
\section{3. The Quantale Morphism as Bridge}
\label{sec:orgd509611}

\subsection{3.1 Why Quantales, Not Just Lattices?}
\label{sec:org137e16e}

The lattice structure alone suffices for fixpoint computation (Knaster-Tarski).
The quantale structure — the sequential composition ⊗ — is what makes the
Galois connections between systems meaningful.

Without ⊗, a Galois connection (α, γ) : L ⇌ M between lattices preserves only
order information. With ⊗, a \textbf{quantale morphism} preserves both order AND
sequential composition:

\begin{verbatim}
α(a ⊗ b) = α(a) ⊗' α(b)     (preserves composition)
α(⊔ᵢ aᵢ) = ⊔'ᵢ α(aᵢ)        (preserves joins)
\end{verbatim}

This is critical for Prologos because the systems don't just compute values —
they compute \emph{sequences of operations}. Substitution application, protocol
continuation, and type-checking steps all have sequential structure. The quantale
morphism guarantees that the sequential structure is preserved across the Galois
bridge.

In System II (sessions), this means: the session protocol's sequential structure
(send, then receive, then end) maps faithfully to the effect ordering's
sequential structure (effect at position 0 before effect at position 1 before
effect at position 2). The morphism doesn't just preserve ordering — it preserves
the \emph{reason} for the ordering.
\subsection{3.2 The Bridge Pattern in Implementation}
\label{sec:org9fa435e}

In Prologos, the Galois connection between domains is implemented as a
\textbf{bridge propagator} — a propagator that watches cells in one domain and writes
to cells in another:

\begin{verbatim}
bridge-propagator : Cell[L] × Cell[M] → Propagator
  When Cell[L] refines from l to l':
    write α(l') to Cell[M]      (α is the abstraction)
  When Cell[M] refines from m to m':
    write γ(m') to Cell[L]      (γ is the concretization)
\end{verbatim}

The monotonicity of α and γ guarantees that the bridge propagator is a valid
monotone propagator — it can participate in the propagator network's fixpoint
computation without violating confluence.

Existing bridges in Prologos:
\begin{itemize}
\item Type ⇌ Session bridge (\texttt{session-propagators.rkt})
\item Type ⇌ Interval bridge (from abstract interpretation work)
\item Session ⇌ Effect Position bridge (proposed in Architecture D)
\end{itemize}

Each bridge is a quantale morphism instantiated as a pair of propagators.
\subsection{3.3 Composition of Galois Connections}
\label{sec:org85a6f89}

Galois connections compose: if (α₁, γ₁) : L ⇌ M and (α₂, γ₂) : M ⇌ N, then
(α₂ ∘ α₁, γ₁ ∘ γ₂) : L ⇌ N is also a Galois connection.

This means the inter-stratum connections compose into a single Galois connection
from the substrate lattice to the barrier lattice. For System III (type checker):

\begin{verbatim}
(α₂ ∘ α₁, γ₁ ∘ γ₂) : T ⇌ A

where α₁ : T* → R (readiness extraction)
      α₂ : R* → A (action construction)
      γ₂ : A → R (action interpretation)
      γ₁ : R → T (resolution commitment)
\end{verbatim}

The \emph{composed} Galois connection maps directly from type cell fixpoints to
resolution actions. This is exactly what the stratified quiescence loop computes
— but the categorical perspective reveals that the intermediate lattice R
(readiness) is not essential; it is a \emph{factorization} of the composed connection
that makes the computation more efficient (readiness can be computed incrementally
by propagators).
\section{4. Formal Statements}
\label{sec:org21a9a52}

We now state the key results precisely.
\subsection{4.1 Definition: Stratifiable Operation}
\label{sec:orgbc321c2}

A function f : L → L on a complete lattice L is \emph{stratifiable with respect to
a monotone endofunctor F : L → L} if the composition f ∘ lfp(F) : L → L is
monotone, where lfp(F) denotes the parametric least fixpoint of F (which is
monotone in the initialization point by Cousot \& Cousot 1979).
\subsection{4.2 Theorem: Generalized Stratification}
\label{sec:org9fb1a80}

Let L be a complete lattice, F : L → L monotone, and f : L → L stratifiable
with respect to F. Define the stratified iteration:

\begin{verbatim}
G = f ∘ lfp(F) : L → L

x₀ = ⊥
xₙ₊₁ = G(xₙ)
\end{verbatim}

Then:
\begin{enumerate}
\item G is monotone (by definition of stratifiability).
\item lfp(G) exists (by Knaster-Tarski).
\item The Kleene chain x₀ ≤ x₁ ≤ \ldots{} converges to lfp(G).
\item lfp(G) is the unique minimal solution to: x = f(lfp(F | x)).
\end{enumerate}
\subsection{4.3 Theorem: Quantale Morphism Preservation}
\label{sec:org69e756b}

Let (Q₁, ≤₁, ⊗₁) and (Q₂, ≤₂, ⊗₂) be quantales and (α, γ) : Q₁ ⇌ Q₂ a
Galois connection that is also a quantale morphism (α preserves ⊗ and ⊔). If
F₁ : Q₁ → Q₁ and F₂ : Q₂ → Q₂ are monotone endofunctors such that
α ∘ F₁ = F₂ ∘ α (the bridge commutes with propagation), then:

\begin{verbatim}
α(lfp(F₁)) ≤ lfp(F₂)
\end{verbatim}

and if additionally γ ∘ F₂ ≤ F₁ ∘ γ, then:

\begin{verbatim}
α(lfp(F₁)) = lfp(F₂)
\end{verbatim}

\emph{This is the soundness theorem for cross-domain bridges.}
It guarantees that fixpoints computed in one domain (e.g., session types) map
correctly to fixpoints in another domain (e.g., effect positions).
\subsection{4.4 Corollary: Confluence Under Non-Overlapping Resolution}
\label{sec:org04f0d80}

For System III (type checker), the barrier operation (trait resolution
commitment) is stratifiable if and only if the instance resolution function is
\emph{deterministic} — i.e., for each (trait, type-args) pair, there is at most one
valid instance.

Under the non-overlapping instance invariant enforced by Prologos (monomorphic
lookup by key, parametric most-specific-wins, same-specificity ties rejected as
HKT-7 ambiguity), the resolution function is deterministic, hence stratifiable,
hence the stratified fixpoint is confluent.

This result was verified empirically by Track 2 Phase 9 (19 confluence tests,
randomized evaluation order, deterministic results).
\section{5. Relationship to Approximation Fixpoint Theory}
\label{sec:org7fdfd5c}

\subsection{5.1 The Connection}
\label{sec:org80f0459}

Approximation Fixpoint Theory (AFT), developed by Denecker, Marek, and
Truszczyński (2000, 2004), provides an algebraic framework for fixpoints of
non-monotone operators using \emph{approximating operators} on bilattices.

The key construction: given a non-monotone operator O : L → L on a complete
lattice L, define an approximating operator A : L² → L² on the bilattice L²
(pairs (x, y) where x ≤ y, ordered by the \emph{precision} ordering ≤\textsubscript{p}), such that
A is monotone with respect to ≤\textsubscript{p}. The fixpoints of A approximate the fixpoints
of O.
\subsection{5.2 Layered Recovery as a Special Case of AFT}
\label{sec:orgfe00435}

Our stratified recovery systems can be seen as a restricted case of AFT where:

\begin{enumerate}
\item The approximating operator A decomposes into a chain of Galois connections
(our inter-stratum functors).
\item The bilattice structure is implicit in the stratification: the lower bound
of the approximation is the fixpoint of the monotone phase (lfp(F\textsubscript{m})), and
the upper bound is determined by the barrier operation.
\item The precision ordering corresponds to the number of strata computed:
more strata = more precise approximation.
\end{enumerate}

However, Layered Recovery is \emph{more structured} than general AFT because:
\begin{itemize}
\item The monotone and non-monotone phases are cleanly separated (not interleaved).
\item The quantale enrichment provides sequential composition, which AFT does not
require.
\item The Galois connections between strata provide soundness guarantees that
general approximating operators do not.
\end{itemize}
\subsection{5.3 What AFT Offers for Future Extensions}
\label{sec:org7dd2d02}

AFT provides semantics for operations that Layered Recovery currently does not
handle well:

\begin{itemize}
\item \textbf{Well-founded semantics}: For self-referential negation (P ← ¬P), AFT assigns
the well-founded model (unknown). Layered Recovery's stratification rejects
this as unstratifiable. AFT provides a graceful degradation.
\item \textbf{Stable models}: For programs with multiple stable fixpoints (non-deterministic
choice), AFT characterizes all stable models. Layered Recovery computes one
fixpoint (the least). AFT could extend the ATMS layer to enumerate alternatives.
\item \textbf{Three-valued fixpoints}: AFT's bilattice naturally supports three-valued logic
(true, false, unknown). This could extend Prologos's type checker to track
constraints whose satisfiability is not yet determined, rather than failing
eagerly.
\end{itemize}
\section{6. Variants and Approaches for Non-Monotonic Domains}
\label{sec:orgefa7a79}

\subsection{6.1 Classification of Non-Monotone Operations}
\label{sec:org416f26a}

Not all non-monotone operations are stratifiable. We classify them:
\subsubsection{Class A: Stratifiable (barrier-compatible)}
\label{sec:orgb870b70}

The operation f : L → L is stratifiable if f ∘ lfp(F) is monotone. This means
adding more information to the substrate's fixpoint never reverses the barrier
operation's output.

\emph{Examples}:
\begin{itemize}
\item Negation-as-failure (when the program is stratifiable)
\item Effect execution (when the effect ordering is derived from a monotone source)
\item Trait resolution commitment (under non-overlapping instances)
\end{itemize}

\emph{Approach}: Direct application of the Layered Recovery Principle. Insert a
barrier between the monotone fixpoint and the non-monotone operation.
\subsubsection{Class B: Locally Stratifiable (barrier-compatible with context restriction)}
\label{sec:org3ad20d0}

The operation f is not globally stratifiable but becomes stratifiable when
restricted to a subset of the lattice. This typically happens when non-monotonicity
arises only for certain value combinations.

\emph{Examples}:
\begin{itemize}
\item Overlapping instance resolution (stratifiable for non-overlapping subsets)
\item Aggregation with non-monotone aggregates (e.g., average — stratifiable when
the set being aggregated is fixed)
\end{itemize}

\emph{Approach}: \emph{Dynamic stratification}. Partition the lattice at runtime into
regions where f is stratifiable, and process each region at a separate barrier.
This is analogous to local stratification in Datalog.
\subsubsection{Class C: Anti-Monotone (requires bilattice or ATMS)}
\label{sec:org7e951e1}

The operation f is anti-monotone: more input information decreases the output.
Classical negation (¬) is the canonical example.

\emph{Examples}:
\begin{itemize}
\item Classical negation (not NAF — actual logical negation)
\item Set complement
\item Constraint \emph{retraction} (removing a previously-added constraint)
\end{itemize}

\emph{Approach}: Use AFT's bilattice construction. Lift the operation to L² where
the approximating operator is monotone. Or use the ATMS: treat the anti-monotone
operation's inputs as hypotheses, compute fixpoints for each hypothesis set,
and select the consistent worldview.
\subsubsection{Class D: Chaotic (no fixpoint guarantee)}
\label{sec:org0f809dc}

The operation has no monotonicity or anti-monotonicity structure. Applying it
repeatedly may oscillate without converging.

\emph{Examples}:
\begin{itemize}
\item Arbitrary function composition with side effects
\item Operations that depend on evaluation order (inherently non-confluent)
\end{itemize}

\emph{Approach}: \emph{Bounded iteration with fuel}. Iterate the stratified chain up to
a fuel limit (as System III does with fuel=100). If convergence is not reached,
report an error. This is not a theoretical guarantee — it is a pragmatic safety
valve. Alternatively, impose ordering constraints to make the operation
deterministic (converting Class D to Class A at the cost of parallelism).
\subsection{6.2 The Adaptation Recipe}
\label{sec:org2688645}

For any new non-monotonic domain:

\begin{enumerate}
\item \textbf{Identify the substrate quantale (L, ≤, ⊗)}: What is the lattice of partial
information? What is sequential composition? What is the join (information merge)?

\item \textbf{Identify the non-monotone operation f}: What operation on L violates
monotonicity? Classify it (A, B, C, or D above).

\item \textbf{For Class A}: Define the monotone propagation phase F\textsubscript{m} and verify
stratifiability. Insert a single barrier.

\item \textbf{For Class B}: Identify the stratifiable partitions. Design dynamic
stratification that detects partition boundaries at runtime.

\item \textbf{For Class C}: Choose between AFT (bilattice approximation) and ATMS
(hypothetical reasoning). AFT is simpler for pure logic; ATMS is better
when the non-monotone operation interacts with choice.

\item \textbf{For Class D}: Apply fuel-bounded iteration. Consider whether ordering
constraints can reduce to Class A.

\item \textbf{If cross-domain interaction is needed}: Define a Galois connection
(α, γ) : L ⇌ M between the new domain and existing domains. Verify that
it is a quantale morphism (preserves ⊗). Implement as a bridge propagator.
\end{enumerate}
\subsection{6.3 Specific Future Applications in Prologos}
\label{sec:org475e38f}

\begin{center}
\begin{tabular}{llll}
Domain & Non-Monotone Operation & Class & Approach\\
\hline
Refinement types & Predicate weakening under subtyping & B & Dynamic stratification by predicate structure\\
Dependent types & Type-level computation with recursion & D (bounded) & Fuel-bounded normalization (already exists)\\
Module system & Re-export shadowing & A & Barrier at module boundary\\
Incremental compilation & Definition invalidation & C & ATMS with assumptions per definition\\
Proof search & Backtracking & C & ATMS (existing)\\
Probabilistic types & Bayesian conditioning (non-monotone on probabilities) & B & Stratify by observation order\\
\end{tabular}
\end{center}
\section{7. Implications and Conclusions}
\label{sec:orgaf133f8}

\subsection{7.1 What the Categorical Structure Tells Us}
\label{sec:org654e302}

The three implementations of Layered Recovery are not \emph{analogous} — they are
\emph{instances of the same mathematical structure}. The common structure is a
stratified endofunctor on a quantale-enriched lattice, with Galois connections
between strata and a non-monotone barrier at exactly one point in the chain.

This is not a coincidence. The CALM theorem (Hellerstein 2011) tells us that
non-monotone operations require coordination, and that monotone operations do
not. The Layered Recovery Principle is the \emph{minimal} coordination strategy: all
monotone reasoning happens coordination-free in the propagator network, and the
single barrier provides the coordination required by the CALM theorem. The
quantale enrichment ensures that sequential structure is preserved across
domain bridges.
\subsection{7.2 What This Means for Implementation}
\label{sec:org113a45b}

For implementers, the categorical structure provides:

\begin{enumerate}
\item \textbf{A checklist}: For each new domain, identify the lattice, the composition,
the non-monotone operation, and its class. The categorical structure tells
you exactly what infrastructure is needed.

\item \textbf{A correctness criterion}: If your Galois connection is a quantale morphism
and your barrier operation is stratifiable, confluence is guaranteed. If any
of these conditions fails, the structure tells you exactly where the problem
is.

\item \textbf{A composition principle}: Galois connections compose. Bridge propagators
compose. New domains can be added to the propagator network without modifying
existing domains — just add a bridge.
\end{enumerate}
\subsection{7.3 What This Means for Theorists}
\label{sec:org4c0648c}

The contribution is the identification of quantale enrichment as the key
structure that distinguishes Layered Recovery from general stratified fixpoint
computation. Lattice-theoretic stratification (Datalog, AFT) handles the
fixpoint side. The quantale morphism handles the \emph{sequential composition} side —
which is essential for programming language semantics where operations have
causal order.

The three systems demonstrate that the same categorical pattern arises
independently in logic programming (unification quantale), protocol verification
(session quantale), and type inference (substitution quantale). This suggests
the pattern is fundamental to any system that combines monotone constraint
propagation with non-monotone decision procedures.
\subsection{7.4 Open Questions}
\label{sec:orgc61c6f5}

\begin{enumerate}
\item \textbf{Is there a 2-categorical structure?} The Galois connections between domains
are 1-cells in a 2-category. Do the bridge propagators form 2-cells? If so,
what additional coherence conditions arise?

\item \textbf{Can the ATMS be characterized as a specific stratification?} Currently, the
ATMS is treated as a separate layer. Can it be unified into the stratified
endofunctor framework — perhaps as an internal hom in the quantale?

\item \textbf{What is the complexity of stratifiability checking?} For System III, we
verify stratifiability by enforcing non-overlapping instances (a syntactic
condition). Is there a general decision procedure for stratifiability of
barrier operations?

\item \textbf{Quantaloid enrichment}: When multiple quantales interact via multiple Galois
connections, the natural home is a \emph{quantaloid} (a category enriched over
quantales). Does the full Prologos propagator network, with all its
cross-domain bridges, form a quantaloid-enriched category?
\end{enumerate}
\section{8. References}
\label{sec:orgfad1baf}

\subsection{8.1 Categorical and Algebraic Foundations}
\label{sec:orgec69fb9}

\begin{itemize}
\item Rosenthal, K.I. (1990). \emph{Quantales and Their Applications}. Longman.
\item Stubbe, I. (2013). "An introduction to quantaloid-enriched categories."
\href{https://www-lmpa.univ-littoral.fr/\~stubbe/PDF/SurveyQCats.pdf}{PDF}
\item Paseka, J. \& Rosický, J. (2000). "Quantales." In \emph{Current Research in
Operational Quantum Logic}.
\href{https://www.sciencedirect.com/science/article/abs/pii/S1570795407050061}{ScienceDirect}
\item Fong, B. \& Spivak, D.I. (2019). \emph{Seven Sketches in Compositionality}. §1.3
Galois Connections.
\href{https://math.libretexts.org/Bookshelves/Applied\_Mathematics/Seven\_Sketches\_in\_Compositionality:\_An\_Invitation\_to\_Applied\_Category\_Theory\_(Fong\_and\_Spivak)/01:\_Generative\_Effects\_-\_Orders\_and\_Adjunctions/1.03:\_Galois\_Connections}{LibreTexts}
\end{itemize}
\subsection{8.2 Effect Systems and Graded Monads}
\label{sec:org0198fb9}

\begin{itemize}
\item Katsumata, S. (2014). "Parametric effect monads and semantics of effect
systems." \emph{POPL 2014}.
\href{https://www.researchgate.net/publication/262368930\_Parametric\_Effect\_Monads\_and\_Semantics\_of\_Effect\_Systems}{ResearchGate}
\item Gordon, C.S. (2021). "Lifting Sequential Effects to Indexed
Computation Types." \emph{ESOP 2021}.
\item Orchard, D. et al. (2020). "Unifying graded and parameterised monads."
\href{https://arxiv.org/pdf/2001.10274}{arXiv}
\item Uustalu, T. (2018). "Graded monads and quantified computational effects."
\href{https://mta.ca/\~rrosebru/FMCS2018/Slides/Uustalu.pdf}{Slides}
\item nLab. "Graded monad."
\href{https://ncatlab.org/nlab/show/graded+monad}{nLab}
\item Barreto, S. \& Milius, S. (2025). "A Category-Theoretic Framework for
Dependent Effect Systems."
\href{https://arxiv.org/html/2601.14846}{arXiv}
\end{itemize}
\subsection{8.3 Fixpoint Theory and Stratification}
\label{sec:org0ebc88e}

\begin{itemize}
\item Cousot, P. \& Cousot, R. (1977). "Abstract interpretation: a unified lattice
model for static analysis." \emph{POPL 1977}.
\item Cousot, P. \& Cousot, R. (2014). "A Galois connection calculus for abstract
interpretation." \emph{POPL 2014}.
\href{https://www.di.ens.fr/\~cousot/publications.www/CousotCousot-POPL14-ACM-p2-3-2014.pdf}{PDF}
\item Denecker, M. et al. (2000, 2004). Approximation Fixpoint Theory.
\href{https://proceedings.kr.org/2021/32/kr2021-0032-heyninck-et-al.pdf}{Recent extension}
\item Heyninck, J. et al. (2025). "A Category-Theoretic Perspective on
Approximation Fixpoint Theory."
\href{https://arxiv.org/pdf/2502.09234}{arXiv}
\href{https://www.cambridge.org/core/journals/theory-and-practice-of-logic-programming/article/categorytheoretic-perspective-on-higherorder-approximation-fixpoint-theory/D3A56D482FE54B7DA41C155021088E39}{Cambridge}
\end{itemize}
\subsection{8.4 CALM Theorem and Distributed Monotonicity}
\label{sec:org8493005}

\begin{itemize}
\item Hellerstein, J.M. (2010, 2019). "Keeping CALM: When Distributed Consistency
is Easy." \emph{CACM}.
\href{https://arxiv.org/abs/1901.01930}{arXiv}
\href{https://cacm.acm.org/research/keeping-calm/}{ACM}
\item Conway, N. et al. (2012). "Logic and Lattices for Distributed Programming."
\href{https://dsf.berkeley.edu/papers/UCB-lattice-tr.pdf}{PDF}
\item Alvaro, P. et al. (2011). "Consistency Analysis in Bloom: a CALM and
Collected Approach." \emph{CIDR 2011}.
\href{https://people.ucsc.edu/\~palvaro/cidr11.pdf}{PDF}
\end{itemize}
\subsection{8.5 Propagator Networks}
\label{sec:org3571a28}

\begin{itemize}
\item Radul, A. \& Sussman, G.J. (2009). "The Art of the Propagator." MIT CSAIL
Tech Report.
\item Radul, A. (2009). \emph{Propagation Networks: A Flexible and Expressive Substrate
for Computation}. PhD thesis, MIT.
\end{itemize}
\subsection{8.6 Session Types}
\label{sec:orga3fdac7}

\begin{itemize}
\item Caires, L. \& Pfenning, F. (2010). "Session Types as Intuitionistic Linear
Propositions." \emph{CONCUR 2010}.
\item Wadler, P. (2012). "Propositions as Sessions." \emph{ICFP 2012}.
\item Atkey, R. (2009). "Parameterised Notions of Computation." \emph{JFP 19(3-4)}.
\end{itemize}
\subsection{8.7 Datalog and Non-Monotonic Reasoning}
\label{sec:org5b3509b}

\begin{itemize}
\item Madsen, M. et al. (2024). "From Datalog to Flix: A Declarative Language for
Fixed Points on Lattices."
\href{https://inst.eecs.berkeley.edu/\~cs294-260/sp24/2024-02-12-flix}{Berkeley}
\item Charalambidis, A. et al. (2025). "The Power of Negation in Higher-Order
Datalog." \emph{TPLP}.
\href{https://www.cambridge.org/core/journals/theory-and-practice-of-logic-programming/article/power-of-negation-in-higherorder-datalog/F29B686E15B1D940647C9C500E335194}{Cambridge}
\item Kuper, L. \& Newton, R.R. (2023). "Breaking the Negative Cycle." \emph{ECOOP 2023}.
\href{https://drops.dagstuhl.de/storage/00lipics/lipics-vol263-ecoop2023/LIPIcs.ECOOP.2023.31/LIPIcs.ECOOP.2023.31.pdf}{PDF}
\end{itemize}
\subsection{8.8 Prologos Internal Documents}
\label{sec:org43d3c2a}

\begin{itemize}
\item \texttt{EFFECTFUL\_COMPUTATION\_ON\_PROPAGATORS.org} — Layered Recovery Principle (canonical)
\item \texttt{2026-03-06\_SESSION\_TYPES\_AS\_EFFECT\_ORDERING.org} — System II design
\item \texttt{2026-02-24\_LOGIC\_ENGINE\_DESIGN.org} — System I design
\item \texttt{2026-03-13\_TRACK2\_REACTIVE\_RESOLUTION\_DESIGN.md} — System III design
\item \texttt{2026-02-27\_1026\_GALOIS\_CONNECTIONS\_ABSTRACT\_INTERPRETATION.md} — Galois connections research
\item \texttt{2026-03-11\_PROPAGATOR\_FIRST\_PIPELINE\_AUDIT.md} — Full system audit
\item \texttt{DESIGN\_PRINCIPLES.org} — Data Orientation principle
\end{itemize}
\end{document}
